% Context from chapters-tex/inverse-problems.tex; for mathematical reference, not compilation.
\section{Regularization of Inverse Problems}

We extend the convex regularization introduced in Chapter \ref{chap-variational} to linear measurement operators.

We consider a bounded linear map $\Phi : \Ss \rightarrow \Hh$, where the signal space $\Ss$ is a Hilbert space or, more generally, a Banach space. The data space $\Hh$ is a Hilbert space.
 
The operator models acquisition: an unknown high-resolution signal $f_0 \in \Ss$ gives rise to the noisy observation
\eq{
	y = \Phi f_0 + w \in \Hh
}
where $w \in \Hh$ represents acquisition noise. Here the noise is deterministic; we assume only that $\norm{w}_\Hh$ is bounded.

An acquisition device records finitely many observations, so applications usually take $\Hh=\RR^P$. The number $P$ may be small relative to the desired reconstruction dimension.
 
For numerical implementation, we discretize the signal space as $\Ss=\RR^N$, where $N$ is the number of grid points. Usually the grid is much finer than the measurement set, so $N \gg P$.
 
Infinite-dimensional function spaces remain useful for modeling the unknown $f_0$ and analyzing recovery, particularly when choosing the signal space $\Ss$.

Direct inversion may be impossible because $\Phi$ is not injective or the data do not lie in its range. Even when an inverse exists on the range, it may have a large norm or be unbounded. Applying it to noisy data amplifies the perturbation: formally, $\Phi^{-1}y=f_0+\Phi^{-1}w$.

We now give a few representative examples of forward operators $\Phi$.

   
\paragraph{Denoising.}

The identity operator $\Phi=\Id_\Ss$ with $\Ss=\Hh$ gives the denoising problem studied in Chapters \ref{chap-denoising} and \ref{chap-variational}.

   
\paragraph{De-blurring and super-resolution.}

For a general operator $\Phi$, recovering $f_0$ requires both inversion and denoising. These goals often conflict because inversion amplifies noise.
 
For example, in deblurring, $\Phi$ is a translation-invariant operator corresponding to low-pass filtering with a kernel $h$
\eql{\label{eq-deblurring-operator}
	\Phi f = f \star h.
}
A periodic model takes $\Ss=\Hh=L^2(\TT^d)$; see Proposition~\ref{prop-ti-convol-l2}.
 
In practice, the blurred signal is sampled on a grid: $\Phi f = \enscond{(f \star h)(x_k)}{0 \leq k< P}$. Figure \ref{fig-ip}, middle, shows the resulting low-resolution image $\Phi f_0$.
 
Inverting such operators is useful for increasing the resolution of digital photographs and videos.


   
\paragraph{Interpolation and inpainting.}

Inpainting reconstructs missing pixels. We model the masking operation by an operator that is diagonal in the spatial domain:
\eql{\label{eq-inp-operator}
	(\Phi f)(x) = \choice{
	0 \qifq x \in \Om,\\
	f(x) \qifq x \notin \Om.
	}
}
Here $\Omega$ is the missing region, either a subset of $[0,1]^d$ in the continuous model or a set of pixel indices in the discrete model.
 
Figure \ref{fig-ip}, right, shows an example of a damaged image $\Phi f_0$.

\myfigure{
\tabtrois{
\image{invpbm}{.3}{ip-sample-image}&
\image{invpbm}{.3}{ip-sample-superresol}&
\image{invpbm}{.3}{ip-sample-inpainting}\\
Original $f_0$ & Low resolution $\Phi f_0$ & Masked $\Phi f_0$
}
}{ 
	Examples of forward operators in inverse problems.  
}{fig-ip}


   
\paragraph{Medical imaging.}

Many medical imaging devices provide indirect access to the signal of interest, with acquisition processes that can be approximated by a linear operator $\Phi$.
 
For a tomographic scanner, the Radon transform models acquisition. The Fourier slice theorem relates these measurements to Fourier data along radial lines.
  
An idealized magnetic resonance imaging (MRI) model also uses partial Fourier measurements:
\eql{\label{eq-partial-fourier-operator}
	\Phi f= \enscond{ \hat f(x) }{ x \in \Om }.
}
Here $\Om$ consists of radial lines for tomography or acquisition trajectories such as spirals for MRI.

Electroencephalography (EEG) and magnetoencephalography (MEG) also infer internal activity from measurements at sensors outside the sources. Abstractly, such a linearized model can be written as $(\Phi f)(s)=\int K(s,x)f(x)\,\mathrm{d}x$, where the kernel depends on the geometry and physical model. Unlike a simple image blur, this kernel need not be translation invariant.

   
\paragraph{Regression for supervised learning.}

Linear supervised learning is closely related to the imaging problems studied here, although it uses different notation. Section~\ref{sec-regression} develops the connection between regression and inverse problems.
 
In statistical learning, the data are $n$ pairs $(x_i,y_i)_{i=1}^n$ with feature vectors $x_i \in \RR^p$. A linear prediction model has the form $y_i = \dotp{\be}{x_i}$, with unknown parameter $\be \in \RR^p$. Placing the vectors $x_i$ in the rows of $X \in \RR^{n \times p}$ gives the approximate system $X \be \approx y$. This corresponds to the inverse problem $\Phi f=y$ under the substitutions $\Phi \mapsto X$ and $f \mapsto \be$, with dimensions $(P,N) \rightarrow (n,p)$.
 
In statistical learning, the model need not be correctly specified, and the design matrix $X$ is often random. Its sampling fluctuations must then be controlled as $n\to\infty$.

Ridge regression estimates the parameter by solvi